\documentclass[a4paper,11pt]{article}

\usepackage[utf8]{inputenc}
\usepackage[T1]{fontenc}
\usepackage[german]{babel}
\usepackage[a4paper,margin=2.5cm]{geometry}
\usepackage{enumitem}
\usepackage{titlesec}
\usepackage{longtable}
\usepackage{setspace}
\usepackage{parskip}

\begin{document}

% ---------------------------
% Deckblatt
% ---------------------------

\begin{titlepage}
    \centering

    \vspace*{3cm}

    {\Huge \textbf{Meetingprotokoll}}\\[0.5cm]
    {\Large Statusbericht zur Projektarbeit}\\[2cm]

    \rule{12cm}{0.5pt}\\[0.5cm]

    {\Large Projektarbeit}\\[0.3cm]
    {\large 29. Mai 2026}

    \rule{12cm}{0.5pt}

    \vfill

    {\large THWS / MINOVA}\\

\end{titlepage}

% ---------------------------
% Inhaltsverzeichnis
% ---------------------------

\tableofcontents
\newpage

% ---------------------------
% Inhalt
% ---------------------------

\section{Allgemeine Informationen}

\begin{tabular}{ll}
\textbf{Datum:} & 29.05.2026 \\
\textbf{Art des Meetings:} & Projektmeeting / Statusbesprechung \\
\textbf{Nächstes Meeting:} & 12.06.2026 um 09:00 Uhr \\
\end{tabular}

\vspace{1cm}

\section{Projektstatus}

Der aktuelle Stand des Projekts wurde gemeinsam mit dem Auftraggeber besprochen. Das entwickelte System wurde erfolgreich vorgestellt und getestet.

Frontend und Backend arbeiten bereits erfolgreich zusammen. Die vom System erzeugten Informationen werden per E-Mail versendet und nach einer Bestätigung oder Ablehnung automatisch an die entsprechenden Seiten weitergeleitet. Die Informationen werden dort korrekt dargestellt. Der aktuelle Stand wurde vom Auftraggeber bestätigt.

\subsection{Aktuell abgeschlossene bzw. bearbeitete Aufgaben}

\begin{itemize}[leftmargin=1.2cm]
\item Funktionsfähige Kommunikation zwischen Frontend und Backend umgesetzt
\item Versand von E-Mails mit relevanten Informationen implementiert
\item Verarbeitung von Bestätigungen und Ablehnungen erfolgreich umgesetzt
\item Automatische Weiterleitung der Informationen auf die entsprechenden Übersichtsseiten implementiert
\item Darstellung der Informationen auf den jeweiligen Seiten umgesetzt
\item Gesamtablauf erfolgreich getestet und vom Auftraggeber bestätigt
\end{itemize}

\vspace{0.5cm}

\section{Besprochene Punkte}

Während des Meetings wurden folgende organisatorische und technische Themen besprochen:

\begin{itemize}[leftmargin=1.2cm]
\item Das Preisfeld aus DISPO soll künftig sowohl in den E-Mails als auch in den Informationsansichten dargestellt werden.

\item Für die Zeitfensterbestätigung soll ein zusätzlicher Hinweistext ergänzt werden, beispielsweise:

\textit{„Bitte bestätigen Sie das Zeitfenster bis spätestens [Datum/Uhrzeit].“}

\item Für den Fall einer ausbleibenden Bestätigung soll ein entsprechender Hinweis angezeigt werden, beispielsweise:

\textit{„Sollten wir innerhalb des angegebenen Zeitraums keine Rückmeldung erhalten, werden wir uns mit Ihnen in Verbindung setzen.“}

\item Die Struktur der Projektdokumentation wurde abgestimmt.

\item Der Projektbericht soll unter anderem Projektaufwand, Arbeitszeiten, Reflexion, Projektplan sowie verwendete Technologien enthalten.

\item Die Ergebnisdokumentation soll die Ausgangssituation, die Projektziele und die erzielten Ergebnisse beschreiben.

\item Alle Teammitglieder sollen ihre geleisteten Arbeitsstunden kontinuierlich dokumentieren. Die Aufwände werden später im Projektbericht und in der Projektdokumentation aufgeführt.

\end{itemize}

\vspace{0.5cm}

\section{Nächste Schritte (nächste 2 Wochen)}

\begin{itemize}[leftmargin=1.2cm]
\item Verbesserung und Überarbeitung der E-Mail-Texte
\item Entwicklung einer Kartenansicht für den Kunden zur Anzeige der Fahrerentfernung
\item Aktualisierung des Projektplans
\item Erstellung eines ersten Entwurfs für die Projektdokumentation
\item Erstellung eines ersten Entwurfs für den Projektbericht
\end{itemize}

\vspace{0.5cm}

\section{Nächste Projektphase}

Die nächste Projektphase konzentriert sich auf:

\begin{itemize}[leftmargin=1.2cm]
\item Erweiterung der Informationsdarstellung für Kunden
\item Entwicklung der Kartenansicht
\item Vorbereitung der Projekt- und Ergebnisdokumentation
\end{itemize}

\vspace{0.5cm}

\section{Nächste Meilensteine}

\begin{itemize}[leftmargin=1.2cm]
\item Integration der erweiterten E-Mail-Texte
\item Implementierung der Kartenansicht
\item Aktualisierung des Projektplans
\item Fertigstellung der ersten Dokumentationsentwürfe
\end{itemize}

\vspace{1cm}

\section{Zusammenfassung}

Das Projekt befindet sich aktuell in einer stabilen Entwicklungsphase. Die Kommunikation zwischen Frontend und Backend funktioniert bereits zuverlässig und die Verarbeitung von Bestätigungen sowie Ablehnungen wurde erfolgreich umgesetzt. Der aktuelle Entwicklungsstand wurde dem Auftraggeber vorgestellt und positiv bewertet.

In den kommenden Wochen liegt der Fokus insbesondere auf der Verbesserung der E-Mail-Kommunikation, der Entwicklung einer Kartenansicht für Kunden sowie der Vorbereitung der Projekt- und Ergebnisdokumentation.

\end{document}
